\section{Background and Related Work}


\subsection{Existing DAO Frameworks}

\subsubsection{Aragon}

Aragon \cite{aragon} pioneered modular DAO frameworks with its Agent-based architecture. However, its approach couples voting and execution more tightly than Azorius, and lacks native support for account abstraction. Aragon's upgrade mechanism relies on proxy patterns that can introduce security complexities.

\subsubsection{Compound Governor}

Compound's Governor contract \cite{compound} established the standard for on-chain governance in DeFi protocols. While elegant, its monolithic design makes it difficult to swap voting strategies or implement complex hierarchical relationships. Gas costs for voting remain prohibitive for smaller token holders.

\subsubsection{Moloch DAO}

Moloch \cite{moloch} introduced the concept of "ragequit" for minority protection, focusing on simplicity and security. However, its fixed governance model (one-member-one-vote with guild shares) limits applicability to diverse organizational structures. The lack of modularity prevents experimentation with alternative voting mechanisms.

\subsubsection{DAOstack}

DAOstack \cite{daostack} explored holographic consensus and reputation-based voting. While innovative, its complexity and gas costs hindered widespread adoption. The framework's tight coupling between components made it difficult to adopt incrementally.

\subsection{Account Abstraction (ERC-4337)}

ERC-4337 \cite{erc4337} introduced a standard for account abstraction without requiring protocol-level changes. Key components include:

\begin{itemize}
    \item \textbf{UserOperations}: Transaction-like objects signed by smart contract accounts
    \item \textbf{Bundlers}: Off-chain services that bundle UserOperations into transactions
    \item \textbf{EntryPoint}: Singleton contract that processes UserOperations
    \item \textbf{Paymasters}: Contracts that sponsor gas fees for UserOperations
\end{itemize}

While ERC-4337 has been adopted for wallet infrastructure, its application to DAO governance systems is novel. Lux DAO represents one of the first production implementations of account abstraction for voting mechanisms.

\subsection{Gnosis Safe and Zodiac}

Gnosis Safe \cite{safe} provides battle-tested multisig infrastructure with a modular architecture. The Zodiac framework \cite{zodiac} extends Safe with standardized module interfaces, enabling composable governance tools. Azorius builds on this foundation, implementing the Zodiac pattern for direct Safe integration.

\subsection{Hats Protocol}

Hats Protocol \cite{hats} provides on-chain organizational structures through a tree-based hierarchy of roles (``hats''). Each hat represents permissions and responsibilities within an organization. Lux DAO's integration enables sophisticated role-based governance with automated management.

\subsection{ERC-6551: Token-Bound Accounts}

ERC-6551 \cite{erc6551} enables any ERC-721 NFT to own assets and interact with contracts. By associating each role (hat) with a token-bound account, Lux DAO enables automated payment streaming to role holders, solving the compensation coordination problem in decentralized organizations.

